\Needspace{6\baselineskip}
\begin{center}\hypertarget{native-249-article177}{}\pdfbookmark[2]{177.}{article-177}177.\end{center}
It follows immediately that the most probable system of values of the quantities $p$, $q$, $r$, $s$ etc. will be the one in which $\Omega$ attains its maximum value, and must therefore be found from the $\nu$ equations $\frac{\mathrm{d}\Omega}{\mathrm{d}p}=0$, $\frac{\mathrm{d}\Omega}{\mathrm{d}q}=0$, $\frac{\mathrm{d}\Omega}{\mathrm{d}r}=0$, $\frac{\mathrm{d}\Omega}{\mathrm{d}s}=0$ etc. On setting $M-V=v$, $M'-V'=v'$, $M''-V''=v''$ etc., and $\frac{\mathrm{d}\varphi\Delta}{\varphi\Delta\,.\,\mathrm{d}\Delta}=\varphi'\Delta$, these equations take the following form:
\[
\begin{gathered}
\frac{\mathrm{d}v}{\mathrm{d}p}\varphi'v+
\frac{\mathrm{d}v'}{\mathrm{d}p}\varphi'v'+
\frac{\mathrm{d}v''}{\mathrm{d}p}\varphi'v''+\text{etc.}=0\\
\frac{\mathrm{d}v}{\mathrm{d}q}\varphi'v+
\frac{\mathrm{d}v'}{\mathrm{d}q}\varphi'v'+
\frac{\mathrm{d}v''}{\mathrm{d}q}\varphi'v''+\text{etc.}=0\\
\frac{\mathrm{d}v}{\mathrm{d}r}\varphi'v+
\frac{\mathrm{d}v'}{\mathrm{d}r}\varphi'v'+
\frac{\mathrm{d}v''}{\mathrm{d}r}\varphi'v''+\text{etc.}=0\\
\frac{\mathrm{d}v}{\mathrm{d}s}\varphi'v+
\frac{\mathrm{d}v'}{\mathrm{d}s}\varphi'v'+
\frac{\mathrm{d}v''}{\mathrm{d}s}\varphi'v''+\text{etc.}=0\\
\text{etc.}
\end{gathered}
\]

A fully determinate solution of the problem can thus be derived by elimination as soon as the nature of the function $\varphi'$ is known. Since it cannot be determined a priori, we shall approach the matter from the other side and inquire upon which \SourceLeaf{250}function, tacitly taken as a basis, the familiar principle whose merit is generally recognized actually rests. It is customary to regard as an axiom the assumption that, if a quantity has been determined by several direct observations made under the same circumstances and with equal care, the arithmetic mean of all the observed values gives the most probable value, if not with absolute rigor, at least approximately, so that it is always safest to adhere to it. Setting $V=V'=V''=\text{etc.}=p$, we must therefore have in general $\varphi'(M-p)+\varphi'(M'-p)+\varphi'(M''-p)+\text{etc.}=0$ if we substitute for $p$ the value $\frac{1}{\mu}(M+M'+M''+\text{etc.})$, whatever positive integer $\mu$ represents. Supposing then that $M'=M''=\text{etc.}=M-\mu N$, we shall have in general, i.\,e. for every positive integer value of $\mu$, $\varphi'(\mu-1)N=(1-\mu)\varphi'(-N)$, from which it is easy to infer that $\frac{\varphi'\Delta}{\Delta}$ must in general be a constant quantity, which we shall denote by $k$. Hence $\log\varphi\Delta=\tfrac12 k\Delta\Delta+\text{Const.}$, or, denoting the base of the hyperbolic logarithms by $e$ and supposing $\text{Const.}=\log\varkappa$,
\[
\varphi\Delta=\varkappa e^{\frac12 k\Delta\Delta}.
\]
Further, it is easy to see that $k$ must necessarily be negative if $\Omega$ is actually to attain a maximum, so we shall set $\tfrac12 k=-hh$; and since, by the elegant theorem first discovered by the illustrious \textsc{Laplace}, the integral $\int e^{-hh\Delta\Delta}\mathrm{d}\Delta$, taken from $\Delta=-\infty$ to $\Delta=+\infty$, becomes $=\frac{\BareRoot{\pi}}{h}\mathpunct{\raisebox{.45ex}{,}}$ (denoting by $\pi$ the semicircumference of a circle whose radius is $1$), our function becomes
\[
\varphi\Delta=\frac{h}{\BareRoot{\pi}}e^{-hh\Delta\Delta}.
\]
